% ============================================================================
% Глава 3: Архитектура и реализация
% ============================================================================

\section{АРХИТЕКТУРА И РЕАЛИЗАЦИЯ}

\subsection{Структура на данните}

Базата данни е организирана в три основни колекции, които отразяват бизнес логиката на платформата за електронно обучение. Дизайнът балансира между влагане на документи (embedded documents) за тясно свързани данни и референции за поддържане на връзки между отделни колекции \cite{banker2011mongodb}.

\subsubsection{Колекция users}

Колекцията съхранява профилите на потребителите на системата - студенти, инструктори и администратори. Структурата включва основни лични данни и предпочитания за персонализация.

\begin{lstlisting}[language=JavaScript,caption=Примерен документ от колекцията users]
{
  "_id": ObjectId("507f1f77bcf86cd799439011"),
  "name": "Иван Петров",
  "email": "ivan.petrov@university.edu",
  "interests": ["Data Science", "Machine Learning"],
  "role": "student",
  "enrollment_date": "2023-09-01T00:00:00Z",
  "is_active": true
}
\end{lstlisting}

\subsubsection{Колекция courses}

Съдържа информация за наличните курсове. Модулите и лекциите са вложени директно в документа, тъй като представляват неделима част от курса и рядко се достъпват независимо.

\begin{lstlisting}[language=JavaScript]
{
  "_id": ObjectId("507f1f77bcf86cd799439012"),
  "title": "Въведение в Data Science",
  "lecturer": "д-р Божидар Иванов",
  "category": "Data Science",
  "price": 299.99,
  "modules": [
    {
      "title": "Модул 1: Основи на данните",
      "lectures": [
        {
          "title": "Какво са данните?",
          "duration": 45,
          "video_url": "https://example.com/lecture1.mp4"
        }
      ]
    }
  ],
  "created_at": "2023-01-15T00:00:00Z",
  "is_active": true
}
\end{lstlisting}

\subsubsection{Колекция progress}

Проследява напредъка на студентите по различните курсове. Използва референции към потребители и курсове, за да избегне дублиране на данни, като същевременно позволява гъвкавост в записването на допълнителни метаданни.

\begin{lstlisting}[language=JavaScript]
{
  "_id": ObjectId("507f1f77bcf86cd799439013"),
  "student_id": ObjectId("507f1f77bcf86cd799439011"),
  "course_id": ObjectId("507f1f77bcf86cd799439012"),
  "status": "in_progress",
  "enrollment_date": "2023-09-15T00:00:00Z",
  "current_module": 1,
  "current_lecture": 2,
  "test_results": [
    {
      "module": 1,
      "score": 85,
      "attempts": 1,
      "date_taken": "2023-10-01T00:00:00Z"
    }
  ],
  "last_activity": "2023-10-20T14:30:00Z",
  "alternative_completion": {
    "type": "project",
    "description": "Изграждане на визуализация на данни",
    "submission_date": "2023-11-01T00:00:00Z",
    "grade": "A"
  }
}
\end{lstlisting}

Тази структура демонстрира гъвкавостта на документо-ориентирания модел - полето \texttt{alternative\_completion} се появява само при студенти, които използват алтернативни методи за завършване, без да се налага добавяне на NULL стойности в релационна база данни \cite{json_schema}.

\subsection{Изпълнявани операции}

\subsubsection{Управление на достъпа (RBAC)}

Системата използва вградения механизъм за контрол на достъпа, базиран на роли (Role-Based Access Control - RBAC) на MongoDB. Сигурността е реализирана на ниво база данни чрез специален конфигурационен скрипт, който дефинира три основни потребителски роли:

\begin{itemize}
    \item \textbf{DataAnalyst}: Роля с права само за четене (read-only) върху всички колекции (\texttt{users}, \texttt{courses}, \texttt{progress}) за целите на анализа на данни.
    \item \textbf{Instructor}: Роля с пълни права (CRUD) върху колекциите \texttt{courses} и \texttt{progress}, както и права за четене на потребителски профили.
    \item \textbf{Student}: Роля с права за четене върху курсовете и ограничени права за редактиране на собствения прогрес.
\end{itemize}

Създаването на ролите се осъществява чрез командата \texttt{createRole} в административната база данни (\texttt{admin}), след което се създават потребители със съответните роли.

\begin{lstlisting}[language=JavaScript,caption=Създаване на потребителски роли и присвояване на привилегии]
// Пример за създаване на роля Instructor
db.command({
  createRole: 'Instructor',
  privileges: [
    { 
      resource: { db: 'edtech_analytics', collection: 'users' }, 
      actions: ['find']
    },
    {
      resource: { db: 'edtech_analytics', collection: 'courses' }, 
      actions: ['find', 'insert', 'update', 'remove']
    }
  ],
  roles: []
});
\end{lstlisting}

Този подход гарантира високо ниво на сигурност, предотвратявайки неоторизиран достъп директно на ниво сървър.

\subsubsection{CRUD операции}

Системата реализира пълния набор от CRUD операции върху трите колекции:

\textbf{Създаване (Create):}
\begin{itemize}
    \item Добавяне на нов потребител с масив от интереси
    \item Създаване на курс с вложени модули и лекции
    \item Регистрация на студент в курс с начални параметри
\end{itemize}

\textbf{Четене (Read):}
\begin{itemize}
\item Филтриране по роля и статус на активност:
\begin{lstlisting}
db.users.find({role: "student", is_active: true})
\end{lstlisting}
\item Търсене по регулярен израз в заглавия на курсове: 
\begin{lstlisting}
db.courses.find({title: {$regex: "Data", $options: "i"}})
\end{lstlisting}
\item Филтриране по диапазон на цени: 
\begin{lstlisting}
db.courses.find({price: {$gte: 200, $lte: 400}})
\end{lstlisting}
\item Комбинирани условия с OR:
\begin{lstlisting}
db.progress.find({$or: [
  {status: "completed"}, 
  {status: "in_progress"}
]})
\end{lstlisting}
\end{itemize}

\textbf{Актуализиране (Update):}
\begin{itemize}
\item Добавяне на нова оценка в масива на резултатите: 
\begin{lstlisting}
db.progress.updateOne({_id: ObjectId("...")}, {$push: 
  {test_results: 
    {module: 2, score: 92, attempts: 1, date_taken: new Date()}
  }
})
\end{lstlisting}
\item Увеличаване на текущата лекция: 
\begin{lstlisting}
db.progress.updateOne({_id: ObjectId("...")}, {$inc: 
  {current_lecture: 1}
})
\end{lstlisting}
\item Актуализиране на множество документи: 
\begin{lstlisting}
db.users.updateMany({role: "student"}, {$set: {is_active: true}})
\end{lstlisting}
\end{itemize}

\textbf{Изтриване (Delete):}
\begin{itemize}
\item Премахване на неактивни профили: 
\begin{lstlisting}
db.users.deleteMany({is_active: false, 
                    last_login: {$lt: new Date("2024-01-01")}})
\end{lstlisting}
\item Изтриване на конкретен запис по ID: 
\begin{lstlisting}
db.progress.deleteOne({_id: ObjectId("...")})
\end{lstlisting}
\end{itemize}

\subsubsection{Агрегиращи операции}

За анализ на данни е реализирана агрегираща заявка за намиране на топ 5 най-активни студенти \cite{mongodb_definitive}:

\begin{lstlisting}[language=JavaScript,caption=Агрегираща заявка за анализ на топ 5 активни студенти]
db.progress.aggregate([
  { $match: { status: 'completed' } },
  {
    $group: {
      _id: '$student_id',
      completedCourses: { $sum: 1 },
      averageScore: { $avg: '$test_results.score' },
      lastActivity: { $max: '$last_activity' }
    }
  },
  { $sort: { completedCourses: -1, averageScore: -1 } },
  { $limit: 5 },
  {
    $lookup: {
      from: 'users',
      localField: '_id',
      foreignField: '_id',
      as: 'student'
    }
  },
  { $unwind: '$student' },
  {
    $project: {
      studentName: '$student.name',
      email: '$student.email',
      completedCourses: 1,
      averageScore: { $round: ['$averageScore', 1] },
      lastActivity: 1
    }
  }
])
\end{lstlisting}

Тази заявка демонстрира използването на операторите \texttt{\$match}, \texttt{\$group}, \texttt{\$sort}, \texttt{\$limit}, \texttt{\$lookup} и \texttt{\$project} за комплексен анализ на данни от две колекции.

\subsection{Тестване и обработка на грешки}

Системата включва автоматизирани unit тестове, които покриват както успешните CRUD операции, така и сценарии с невалидни заявки за проверка на обработката на грешки. По-долу са представени логове от тестовете за невалидни и неуспешни заявки:

\textbf{1. Заявка към несъществуващ маршрут (404 Not Found):}
\begin{lstlisting}[language=bash,caption=Лог от заявка към несъществуващ маршрут]
GET /nonexistent
Response Status: 404
Response Body:
{
  "error": "Route not found"
}
\end{lstlisting}

\textbf{2. Заявка с невалиден формат на идентификатора (500 Internal Server Error):}
\begin{lstlisting}[language=bash,caption=Лог от заявка с невалиден ObjectId]
PUT /users/invalid-id
Response Status: 500
Response Body:
{
  "error": "input must be a 24 character hex string, 12 byte Uint8Array, or
an integer"
}
\end{lstlisting}

\textbf{3. Заявка за изтриване на несъществуващ потребител (404 Not Found):}
\begin{lstlisting}[language=bash,caption=Лог от заявка за триене на несъществуващ запис]
DELETE /users/507f1f77bcf86cd799439011
Response Status: 404
Response Body:
{
  "message": "User not found"
}
\end{lstlisting}

Тези тестове гарантират стабилността на приложението при неправилно подадени входни данни.

Архитектурата на системата осигурява ефективно управление на образователни данни с висока степен на гъвкавост и производителност, като същевременно поддържа целостта и консистентността на информацията \cite{martin2017clean}.
